\chapter{Performance Measurement for Time-Series Databases}

Probably the most important aspects of a database are its reliability,
security as well as crash-safety. However, as most databases already provide a
certain level of these features, performance often becomes the deciding factor
when choosing a database for a specific use case. Important performance
indicators for time-series databases include ingestion rate, query response
time, compression efficiency, and scalability.

\section{Key Performance Indicators (KPIs)}

Unlike traditional transaction-heavy databases, TSDBs are evaluated based on
their ability to handle massive streams of time-stamped data and complex
aggregations. The primary KPIs include~\cite{TSBS:01}:

\begin{itemize}
    \item \textbf{Ingestion Rate:} The volume of data points inserted per second. High ingestion rates (measured in metrics or rows per second) are critical for reducing visibility latency in real-time processing applications.
    \item \textbf{Query Response Time:} Real-time analytics and alerting demand
          sub-second query latency. Response times are evaluated across various patterns,
          including recent data point lookups, historical scans, and windowed
          aggregations~\cite{RTABench:01}.
    \item \textbf{Compression Efficiency:} Due to the high volume of historical data, efficient compression algorithms (like delta or run-length encoding) are essential to minimize storage costs and optimize I/O, balancing compression ratios with computational overhead~\cite{Compression:01}.
    \item \textbf{Scalability:} To handle growth from millions to billions of data points, TSDBs employ time-based partitioning and sharding. Scalability measures performance consistency as datasets grow (data scalability) and the speedup achieved by adding hardware nodes, theoretically bounded by Amdahl's Law~\cite{Amdahl:01}.
\end{itemize}

\section{Benchmarking Tools and Methodologies}

Standardized tools ensure objective, reproducible performance comparisons
across different databases by employing deterministic data and query
generation.

\subsection{Time Series Benchmark Suite (TSBS)}

The TSBS is the de facto standard for benchmarking TSDBs (e.g., TimescaleDB,
InfluxDB, MongoDB). It measures both bulk load performance and query execution
latencies using predefined scenarios~\cite{TSBS:01}:

\begin{itemize}
    \item \textbf{DevOps Scenario:} Simulates data center monitoring by generating infrastructure metrics (CPU, memory, disk) and diverse metadata tags for a configurable number of hosts. Queries evaluate aggregations and multi-dimensional filtering.
    \item \textbf{IoT Scenario:} Models fleet management diagnostics, introducing real-world complexities such as temporary batch ingestion, out-of-order data, and missing points, while blending real-time monitoring with predictive analytics.
\end{itemize}

\subsection{Real-Time Analytics Benchmarking (RTABench)}

Modern applications often require capabilities beyond those tested in
traditional TSBS workloads, particularly involving complex state computations
and materialized views. Real-time analytics workloads typically require:
\begin{enumerate}
    \item Fast multi-table joins on fresh event data and normalized metadata.
    \item Selective filtering based on specific entity IDs and narrow time windows.
    \item Sub-second aggregations heavily reliant on incrementally updated materialized
          views.
\end{enumerate}

To accurately measure these distinct query patterns, RTABench supplements
standard tools by evaluating operations over an entity-normalized data model,
demonstrating the vital trade-offs between flexibility and query
optimization~\cite{RTABench:01}. Meaningful benchmark interpretation ultimately
requires mapping test workloads directly to production use cases while
accounting for environmental consistency and proper, fair database
configuration.
